% =========================================================================
% TEORIA DO FLUXO DE CUSTO MÍNIMO E DUALIDADE
% =========================================================================
\chapter{Teoria do Fluxo de Custo Mínimo e Dualidade}\label{sec:mcf_teoria}
\label{sec:mcf}

O problema do Fluxo de Custo Mínimo (\textit{Minimum Cost Flow} --- MCF) generaliza o fluxo máximo associando a cada arco, além da capacidade, um custo unitário a ser minimizado~\cite{ahuja1993network}.

Formalmente, uma rede de custo mínimo é definida sobre um digrafo $D = (V, A)$ onde cada arco $(u,v) \in A$ possui:
\begin{itemize}
    \item \textbf{Capacidade:} $c(u,v) \ge 0$, limite superior de fluxo no arco.
    \item \textbf{Custo unitário:} $w(u,v) \in \mathbb{R}$, custo incorrido por unidade de fluxo no arco.
\end{itemize}

Cada vértice $v \in V$ possui demanda líquida $b(v) \in \mathbb{R}$:
\begin{itemize}
    \item $b(v) > 0$: vértice \textbf{fonte}, fornecendo $b(v)$ unidades de fluxo;
    \item $b(v) < 0$: vértice \textbf{sumidouro}, consumindo $|b(v)|$ unidades;
    \item $b(v) = 0$: vértice de \textbf{transbordo}, com conservação estrita de fluxo.
\end{itemize}

A condição de equilíbrio global $\sum_{v \in V} b(v) = 0$ é necessária para a viabilidade.

\begin{definicao}[Custo Total de um Fluxo][def:custo_fluxo]
Dada a rede $D = (V, A, c, w)$ e um fluxo viável $f$, o \textbf{custo total} $\cost(f)$ é:
\begin{equation}
    \cost(f) \coloneqq \sum_{(u,v) \in A} w(u,v) \cdot f(u,v)
\end{equation}
\end{definicao}

O Problema do Fluxo de Custo Mínimo consiste em:
\begin{align}
    \min_{f} \quad & \sum_{(u,v) \in A} w(u,v) \cdot f(u,v) \label{eq:mcf_obj} \\[4pt]
    \sa \quad & \sum_{(v,w) \in A} f(v,w) - \sum_{(u,v) \in A} f(u,v) = b(v), \quad \forall v \in V \label{eq:mcf_balanco} \\[4pt]
    & 0 \le f(u,v) \le c(u,v), \quad \forall (u,v) \in A \label{eq:mcf_cap}
\end{align}

Esta formulação abrange como casos particulares os problemas de transporte, atribuição, caminho mais curto e fluxo máximo~\cite{cook1998combinatorial, korte2018combinatorial}.

\begin{figure}[!ht]
\centering
\begin{tikzpicture}[thick]
    \node[source node] (s1) at (0, 2.0) {$v_1$};
    \node[source node] (s2) at (0, -2.0) {$v_2$};
    \node[trans node] (m1) at (4.0, 2.0) {$v_3$};
    \node[trans node] (m2) at (4.0, -2.0) {$v_4$};
    \node[sink node] (t1) at (8.0, 2.0) {$v_5$};
    \node[sink node] (t2) at (8.0, -2.0) {$v_6$};

    \node[font=\footnotesize, above=3pt] at (s1.north) {$b=+5$};
    \node[font=\footnotesize, below=3pt] at (s2.south) {$b=+3$};
    \node[font=\footnotesize, above=3pt] at (m1.north) {$b=0$};
    \node[font=\footnotesize, below=3pt] at (m2.south) {$b=0$};
    \node[font=\footnotesize, above=3pt] at (t1.north) {$b=-4$};
    \node[font=\footnotesize, below=3pt] at (t2.south) {$b=-4$};

    \draw[flow edge] (s1) -- node[edge lbl above] {$w\!=\!2, c\!=\!4$} (m1);
    \draw[flow edge] (s1) -- node[edge lbl above, pos=0.25] {$w\!=\!5, c\!=\!3$} (m2);
    \draw[flow edge] (s2) -- node[edge lbl below, pos=0.25] {$w\!=\!3, c\!=\!4$} (m1);
    \draw[flow edge] (s2) -- node[edge lbl below] {$w\!=\!1, c\!=\!5$} (m2);
    \draw[flow edge] (m1) -- node[edge lbl above] {$w\!=\!4, c\!=\!5$} (t1);
    \draw[flow edge] (m1) -- node[edge lbl, right=3pt] {$w\!=\!2, c\!=\!2$} (m2);
    \draw[flow edge] (m2) -- node[edge lbl below, pos=0.25] {$w\!=\!6, c\!=\!3$} (t1);
    \draw[flow edge] (m2) -- node[edge lbl below] {$w\!=\!3, c\!=\!6$} (t2);
\end{tikzpicture}
\caption{Rede de custo mínimo com demandas $b(v)$, custos unitários $w$ e capacidades $c$. Condição de equilíbrio: $\sum b(v) = 5 + 3 - 4 - 4 = 0$.}
\label{fig:mcf_rede}
\end{figure}


\section{Grafo Residual com Custos}

Dado um fluxo viável $f$, o digrafo residual $D_f = (V, A_f)$ é definido com capacidades e custos residuais associados~\cite{ahuja1993network}:
\begin{itemize}
    \item Se $f(u,v) < c(u,v)$: o arco direto $(u,v) \in A_f$ possui capacidade residual $c_f(u,v) = c(u,v) - f(u,v)$ e custo $w_f(u,v) = w(u,v)$.
    \item Se $f(u,v) > 0$: o arco reverso $(v,u) \in A_f$ possui capacidade residual $c_f(v,u) = f(u,v)$ e custo $w_f(v,u) = -w(u,v)$.
\end{itemize}

\begin{figure}[!ht]
\centering
\begin{tikzpicture}[thick]
\begin{scope}[xshift=0cm]
  \node[source node] (s) at (0,1.5) {$s$};
  \node[flow node] (a) at (3.2,3.2) {$a$};
  \node[flow node] (b) at (3.2,-0.2) {$b$};
  \node[sink node] (t) at (6.4,1.5) {$t$};
  \draw[flow edge] (s) -- node[edge lbl above] {$f\!=\!2, w\!=\!3, c\!=\!4$} (a);
  \draw[flow edge] (s) -- node[edge lbl below] {$f\!=\!1, w\!=\!1, c\!=\!3$} (b);
  \draw[flow edge] (a) -- node[edge lbl above] {$f\!=\!2, w\!=\!2, c\!=\!3$} (t);
  \draw[flow edge] (b) -- node[edge lbl below] {$f\!=\!1, w\!=\!5, c\!=\!2$} (t);
  \node[font=\footnotesize, below] at (3.2,-1.0) {\textit{(a) Rede original com fluxo}};
\end{scope}

\node at (7.8,1.5) {\Huge $\Rightarrow$};

\begin{scope}[xshift=9.2cm]
  \node[source node] (s2) at (0,1.5) {$s$};
  \node[flow node] (a2) at (3.2,3.2) {$a$};
  \node[flow node] (b2) at (3.2,-0.2) {$b$};
  \node[sink node] (t2) at (6.4,1.5) {$t$};

  \draw[flow edge,blue] (s2) -- node[edge lbl below, text=blue!80!black] {$c_f\!=\!2, w\!=\!3$} (a2);
  \draw[flow edge,blue] (s2) -- node[edge lbl above, text=blue!80!black] {$c_f\!=\!2, w\!=\!1$} (b2);
  \draw[flow edge,blue] (a2) -- node[edge lbl below, text=blue!80!black] {$c_f\!=\!1, w\!=\!2$} (t2);
  \draw[flow edge,blue] (b2) -- node[edge lbl above, text=blue!80!black] {$c_f\!=\!1, w\!=\!5$} (t2);

  \draw[rev edge,bend right=32] (a2) to node[edge lbl above, text=red!80!black] {$c_f\!=\!2, w\!=\!{-3}$} (s2);
  \draw[rev edge,bend left=32] (b2) to node[edge lbl below, text=red!80!black] {$c_f\!=\!1, w\!=\!{-1}$} (s2);
  \draw[rev edge,bend right=32] (t2) to node[edge lbl above, text=red!80!black] {$c_f\!=\!2, w\!=\!{-2}$} (a2);
  \draw[rev edge,bend left=32] (t2) to node[edge lbl below, text=red!80!black] {$c_f\!=\!1, w\!=\!{-5}$} (b2);
  \node[font=\footnotesize, below] at (3.2,-1.0) {\textit{(b) Digrafo residual $D_f$ com custos}};
\end{scope}
\end{tikzpicture}
\caption{Construção do grafo residual com custos: arcos diretos possuem custo $w$ e arcos reversos possuem custo $-w$.}
\label{fig:custos_residuais}
\end{figure}


\section{Condição de Otimalidade: Ciclos de Custo Negativo}\label{sec:condicao_klein}

\begin{teorema}[Condição de Otimalidade por Ciclos Negativos --- Klein, 1967][teo:ciclos_negativos]
Um fluxo viável $f$ é de custo mínimo se, e somente se, o digrafo residual $D_f$ não contém ciclo direcionado de custo total negativo~\cite{klein1967primal}.
\end{teorema}

\intuicao{Circulação Residual} A diferença $g = f^* - f$ entre dois fluxos viáveis com o mesmo vetor de demandas $b$ possui divergência nula em cada vértice, definindo uma circulação pura. Decomposta em ciclos simples, a diferença de custo $\cost(f^*) - \cost(f)$ expressa-se como a soma dos custos desses ciclos em $D_f$. Se $D_f$ não contém ciclos negativos, nenhuma alteração em ciclo fechado pode reduzir o custo total.

\begin{demonstracao}
$(\Rightarrow)$ Suponha que $f$ é ótimo e existe um ciclo direcionado simples $C$ em $D_f$ com custo total $w(C) < 0$. Seja $\delta = \min_{a \in C} c_f(a) > 0$. Enviando $\delta$ unidades ao longo de $C$, obtém-se um fluxo viável $f'$ com custo:
$$\cost(f') = \cost(f) + \delta \cdot w(C) < \cost(f)$$
contradizendo a otimalidade de $f$.

\medskip
\noindent $(\Leftarrow)$ Suponha que $D_f$ não contém ciclos de custo negativo. Seja $f^*$ um fluxo viável ótimo arbitrário. A função $g = f^* - f$ satisfaz a conservação em todos os vértices e induz uma circulação viável em $D_f$. Pelo Teorema~\ref{teo:decomposicao}, $g$ decompõe-se em ciclos direcionados simples $C_1, \ldots, C_k$ em $D_f$ com $\delta_i > 0$:
$$\cost(f^*) - \cost(f) = \sum_{i=1}^k \delta_i \cdot w(C_i)$$
Como $w(C_i) \ge 0$ para todo $i$, tem-se $\cost(f^*) \ge \cost(f)$. Sendo $f^*$ ótimo, decorre $\cost(f) = \cost(f^*)$, provando a otimalidade de $f$.
\end{demonstracao}

A condição de Klein motiva duas abordagens algorítmicas:
\begin{itemize}
    \item \textbf{Primal (Cycle Canceling):} Inicia com um fluxo viável e cancela iterativamente ciclos negativos em $D_f$.
    \item \textbf{Dual (Successive Shortest Path):} Inicia com fluxo nulo e envia fluxo estritamente por caminhos de custo mínimo, preservando a ausência de ciclos negativos.
\end{itemize}

\begin{figure}[!ht]
\centering
\begin{tikzpicture}[thick]
\begin{scope}[xshift=0cm]
  \tikzsubcaption{(1.75,2.6)}{(a) Ciclo negativo em $D_f$}
  \node[flow node] (a) at (0,1.8) {$a$};
  \node[flow node] (b) at (3.5,1.8) {$b$};
  \node[flow node] (c) at (1.75,0) {$c$};
  \draw[sat edge] (a) -- node[edge lbl above] {$w\!=\!2$} (b);
  \draw[sat edge] (b) -- node[edge lbl above] {$w\!=\!{-4}$} (c);
  \draw[sat edge] (c) -- node[edge lbl above] {$w\!=\!1$} (a);
  \node[font=\footnotesize,below=2pt] at (1.75,-0.7) {$w(C) = 2 + (-4) + 1 = -1 < 0$};
\end{scope}

\node at (5.75,1.0) {\Huge $\Rightarrow$};

\begin{scope}[xshift=7.5cm]
  \tikzsubcaption{(1.75,2.6)}{(b) Após cancelamento}
  \node[flow node] (a2) at (0,1.8) {$a$};
  \node[flow node] (b2) at (3.5,1.8) {$b$};
  \node[flow node] (c2) at (1.75,0) {$c$};
  \draw[tree edge] (a2) -- node[edge lbl above] {$+\delta$} (b2);
  \draw[tree edge] (b2) -- node[edge lbl above] {$+\delta$} (c2);
  \draw[tree edge] (c2) -- node[edge lbl above] {$+\delta$} (a2);
  \node[font=\footnotesize,below=2pt] at (1.75,-0.7) {Custo reduz em $\delta \cdot |w(C)|$};
\end{scope}
\end{tikzpicture}
\caption{(a)~Ciclo negativo $C$ em $D_f$ com $w(C) = -1$. (b)~Cancelamento com aumento $\delta$, reduzindo o custo total.}
\label{fig:ciclo_negativo}
\end{figure}


\section{Custos Reduzidos e Potenciais Nodais}

Dado um vetor de potenciais nodais $\pi: V \rightarrow \mathbb{R}$, o \textbf{custo reduzido} de um arco $(u,v)$ é:
\begin{equation}
    \bar{w}_\pi(u,v) = w(u,v) - \pi(u) + \pi(v)
\end{equation}

Ao longo de qualquer ciclo direcionado $C$, a soma dos custos reduzidos coincide com a dos custos originais, pois os potenciais se anulam telescopicamente:
\begin{equation}
    \sum_{(u,v) \in C} \bar{w}_\pi(u,v) = \sum_{(u,v) \in C} w(u,v) - \sum_{(u,v) \in C} [\pi(u) - \pi(v)] = \sum_{(u,v) \in C} w(u,v)
\end{equation}

\begin{teorema}[Equivalência dos Custos Reduzidos][teo:custos_reduzidos]
Um fluxo viável $f$ é de custo mínimo se, e somente se, existe $\pi: V \to \mathbb{R}$ tal que:
\begin{equation}
    \bar{w}_\pi(u,v) = w_f(u,v) - \pi(u) + \pi(v) \ge 0, \quad \forall\, (u,v) \in A_f
\end{equation}
\end{teorema}

\begin{demonstracao}
$(\Leftarrow)$ Se $\bar{w}_\pi(u,v) \ge 0$ em $A_f$, para qualquer ciclo simples $C$ em $D_f$ tem-se $w(C) = \sum_C \bar{w}_\pi(u,v) \ge 0$. Pelo Teorema~\ref{teo:ciclos_negativos}, $f$ é de custo mínimo.

\medskip
\noindent $(\Rightarrow)$ Se $f$ é ótimo, $D_f$ não possui ciclos negativos. Adicionando uma fonte virtual $s_0$ com arcos de custo zero para todos os vértices, as distâncias mínimas $\dist(s_0, v)$ em relação a $w_f$ são finitas. Definindo $\pi(v) = \dist(s_0, v)$, a desigualdade triangular $\dist(s_0, v) \le \dist(s_0, u) + w_f(u,v)$ implica $\bar{w}_\pi(u,v) \ge 0$.
\end{demonstracao}

A garantia de custos reduzidos não negativos viabiliza o uso de Dijkstra com fila de prioridades no lugar de Bellman-Ford no algoritmo de caminhos aumentantes sucessivos.

\begin{figure}[!ht]
\centering
\begin{tikzpicture}[thick]
    \node[source node] (s) at (0, 0) {$s$};
    \node[flow node] (a) at (3.5, 2.0) {$a$};
    \node[flow node] (b) at (3.5, -2.0) {$b$};
    \node[sink node] (t) at (7, 0) {$t$};

    \node[font=\footnotesize, above=3pt, blue] at (s.north) {$\pi\!=\!0$};
    \node[font=\footnotesize, above=3pt, blue] at (a.north) {$\pi\!=\!2$};
    \node[font=\footnotesize, below=3pt, blue] at (b.south) {$\pi\!=\!1$};
    \node[font=\footnotesize, above=3pt, blue] at (t.north) {$\pi\!=\!5$};

    \draw[flow edge] (s) -- node[edge lbl above] {$w\!=\!3$} node[edge lbl below, text=red!80!black] {$\bar{w}\!=\!5$} (a);
    \draw[flow edge] (s) -- node[edge lbl below] {$w\!=\!2$} node[edge lbl above, text=red!80!black] {$\bar{w}\!=\!3$} (b);
    \draw[flow edge] (a) -- node[edge lbl above] {$w\!=\!4$} node[edge lbl below, text=red!80!black] {$\bar{w}\!=\!7$} (t);
    \draw[flow edge] (b) -- node[edge lbl below] {$w\!=\!6$} node[edge lbl above, text=red!80!black] {$\bar{w}\!=\!10$} (t);
    \draw[flow edge] (a) -- node[edge lbl right] {$w\!=\!1$} node[edge lbl left, text=red!80!black] {$\bar{w}\!=\!0$} (b);
\end{tikzpicture}
\caption{Potenciais nodais $\pi$ (azul) e custos reduzidos $\bar{w} \ge 0$ (vermelho), atestando a otimalidade do fluxo.}
\label{fig:custos_reduzidos}
\end{figure}


\section{Complementaridade de Folgas}

\begin{teorema}[Folgas Complementares para MCF][teo:folga_complementar]
Um fluxo viável $f$ e um vetor de potenciais nodais $\pi$ são ótimos se, e somente se, para todo $(u,v) \in A$:
\begin{enumerate}
    \item Se $\bar{w}_\pi(u,v) > 0$, então $f(u,v) = 0$.
    \item Se $\bar{w}_\pi(u,v) < 0$, então $f(u,v) = c(u,v)$.
    \item Se $0 < f(u,v) < c(u,v)$, então $\bar{w}_\pi(u,v) = 0$.
\end{enumerate}
\end{teorema}

\begin{demonstracao}
As três condições são equivalentes a $\bar{w}_\pi(a) \ge 0$ para todo arco residual $a \in A_f$:
\begin{enumerate}
    \item Se $f(u,v) = 0$, apenas o arco direto existe em $D_f$, exigindo $\bar{w}_\pi(u,v) \ge 0$.
    \item Se $f(u,v) = c(u,v)$, apenas o arco reverso existe em $D_f$, exigindo $\bar{w}_\pi(v,u) = -\bar{w}_\pi(u,v) \ge 0 \implies \bar{w}_\pi(u,v) \le 0$.
    \item Se $0 < f(u,v) < c(u,v)$, ambos os arcos existem em $D_f$, impondo $\bar{w}_\pi(u,v) \ge 0$ e $-\bar{w}_\pi(u,v) \ge 0$, logo $\bar{w}_\pi(u,v) = 0$.
\end{enumerate}
Pelo Teorema~\ref{teo:custos_reduzidos}, essas condições são necessárias e suficientes para otimalidade.
\end{demonstracao}

No algoritmo Network Simplex, o critério de entrada na base identifica arcos fora da base que violam as condições de folga complementar.


\section{Arquitetura Orientada a Objetos da Família CostNetwork}

Os resolvedores de fluxo de custo mínimo herdam da classe-base abstrata \lstinline{CostNetwork} (Código~\ref{lst:costnetwork}, Apêndice~\ref{ap:costnetwork}).

A estrutura de arestas incorpora o custo unitário $w(u,v)$:
\begin{lstlisting}[language=C++, numbers=none, basicstyle=\ttfamily\footnotesize, frame=none, backgroundcolor=\color{white}]
struct Edge {
    Size from, to;
    Long capacity, cost, flow;
};
\end{lstlisting}

A inserção \lstinline{add_edge(from, to, capacity, cost)} constrói em par o arco direto e o reverso:
\begin{itemize}
    \item Arco direto $(u,v)$ com capacidade $c$, custo $w$ e fluxo inicial $0$ (índice $2k$).
    \item Arco reverso $(v,u)$ com capacidade inicial $0$, custo $-w$ e fluxo inicial $0$ (índice $2k+1$).
\end{itemize}

O acesso ao arco complementar opera em $\mathcal{O}(1)$ via \lstinline{id ^ 1ULL}. A interface pública define o método virtual puro \lstinline{compute_min_cost_max_flow(Size source, Size sink)} e métodos \lstinline{make} e \lstinline{clone}.

\paragraph{Antissimetria de Custos em Grafos Não Direcionados.} Em \lstinline{CostNetwork}, conexões bidirecionais exigem a inserção de dois pares independentes:
\begin{itemize}
    \item \lstinline{add_edge(u, v, c, w)}: canal $u \to v$ com reverso de custo $-w$ e capacidade inicial 0.
    \item \lstinline{add_edge(v, u, c, w)}: canal $v \to u$ com reverso de custo $-w$ e capacidade inicial 0.
\end{itemize}
Essa segregação decorre da antissimetria $w_f(v, u) = -w_f(u, v)$, garantindo que o arco reverso represente exclusivamente o cancelamento de fluxo despachado.


